% Generated by thesis/ch7/exp2/make_assets.py -- do not edit by hand.
\begin{table}[tbp]
  \centering
  \small
  \textit{(a) Full against linearised networks}\par\smallskip
  \begin{tabular}{@{}lccc@{}}
    \toprule
    Mean over the ten networks & Full network & Linearised & Gain \\
    \midrule
    Train AUC (242 patients) & 0.956 & 0.927 & $+$0.029 \\
    Test AUC (61 patients) & 0.893 & 0.891 & $+$0.002 \\
    \bottomrule
  \end{tabular}
  \par\bigskip
  \textit{(b) Every non-linear unit and part}\par\smallskip
  \begin{adjustbox}{max width=\textwidth}
  \begin{tabular}{@{}lcccccc@{}}
    \toprule
    & Total & \makecell{Test gain\\CI $>0$} & \makecell{No measurable\\test effect} & Recurring & \makecell{Seed-\\specific} & \makecell{Small support\\($n<30$)} \\
    \midrule
    Non-linear units & 62 & 1 & 61 & 19 & 28 & 16 \\
    Non-linear parts & 45 & 1 & 44 & 9 & 27 & 6 \\
    \bottomrule
  \end{tabular}
  \end{adjustbox}
  \caption{Heart-100 audit: how much of the ten networks' non-linear structure shows on new patients. (a) Each network against itself with every non-linear unit replaced by its best straight line on the training rows; 93\% of the mean training gain does not carry over to the test set, the linearised network does as well or better on test in 5 of 10, and no network has a test-gain interval above zero. (b) A unit or part \emph{generalises} when its own test gain has a 95\% interval above zero; \emph{recurring} means another seed's network reproduces its non-linear residual. Units and parts too close to linear to match (15 and 9) are counted in neither. Binary one-step units are treated as linear. Exploratory, not pre-registered; test numbers were never used to choose.}
  \label{tab:exp2-heart-audit}
\end{table}
